\documentclass[../../main.tex]{subfiles}
\graphicspath{{../../image/}{../../figure/}}
\begin{document}

\section{自相似平面无旋等熵气流}\label{section:ch18-自相似无旋等熵气流}

在平面无旋等熵气流中，速度 $v$ 由速度势 $\phi$ 经\eqref{align:ch03-势流}导出，控制方程为连续性方程\eqref{align:ch03-等熵Euler}$_1$ 与 Bernoulli 方程\eqref{align:ch03-Bernoulli}的二维形式。由于其伸缩性质，这些方程容许形如
\begin{align}
\rho=\rho(\xi),\qquad v=v(\xi),\qquad \phi=t\psi(\xi),\qquad \xi=\frac1t x
\label{align:ch18-1831}
\end{align}
的自相似解。

事实上，在此前提下，以 $\nabla$ 表示关于 $\xi$ 变量的梯度算子，取 $g=0$，则\eqref{align:ch03-势流}、(3.3.36)$_1$ 与 (3.3.49) 成为
\begin{align}
v=\nabla\psi,
\label{align:ch18-1832}
\end{align}
\begin{align}
-\xi\cdot\nabla\rho+\nabla\cdot(\rho\nabla\psi)=0,
\label{align:ch18-1833}
\end{align}
\begin{align}
\psi-\xi\cdot\nabla\psi+\frac12|\nabla\psi|^2+h(\rho)=h(\rho_0).
\label{align:ch18-1834}
\end{align}

在 $\mathbb{R}^2$ 的任一开子集 $\Omega$ 上，若 $\rho$ 取常数值，则\eqref{align:ch18-1833}退化为 $\Delta\psi=0$。此时对\eqref{align:ch18-1834}作用 Laplace 算子得 $|\nabla v|^2=0$，即速度也为常数。沿用前面的术语，这样的 $\Omega$ 称为流动的常状态。

我们可将\eqref{align:ch18-1833}与\eqref{align:ch18-1834}写成更优雅的形式
\begin{align}
\nabla\cdot(\rho\nabla\chi)+2\rho=0,
\label{align:ch18-1835}
\end{align}
\begin{align}
\chi+\frac12|\nabla\chi|^2+h(\rho)=h(\rho_0),
\label{align:ch18-1836}
\end{align}
其中
\begin{align}
\chi(\xi)=\psi(\xi)-\frac12|\xi|^2,
\label{align:ch18-1837}
\end{align}
称为伪势（pseudopotential）。相应地，
\begin{align}
\nabla\chi=v-\xi
\label{align:ch18-1838}
\end{align}
称为伪速度（pseudovelocity），由它生成的“流”称为伪流（pseudoflow）。

特别地，在速度为常值（设为 $\bar v$）的常状态 $\Omega$ 上，伪速度为 $\bar v-\xi$，伪势为
\begin{align}
\chi(\xi)=\bar v\cdot\xi-\frac12|\xi|^2+\text{常数}.
\label{align:ch18-1839}
\end{align}
由\eqref{align:ch18-1834}，在略微滥用记号的意义下，
\begin{align}
\rho=h^{-1}\!\left(h(\rho_0)-\chi-\frac12|\nabla\chi|^2\right)=\rho\bigl(|\nabla\chi|^2,\chi,\rho_0\bigr).
\label{align:ch18-18310}
\end{align}
于是我们可用 $\chi$ 表示声速：
\begin{align}
c^2=p'\bigl(\rho(|\nabla\chi|^2,\chi,\rho_0)\bigr)=c^2\bigl(|\nabla\chi|^2,\chi,\rho_0\bigr).
\label{align:ch18-18311}
\end{align}
将\eqref{align:ch18-18310}中的 $\rho$ 代入\eqref{align:ch18-1835}，导出 $\chi$ 的一个非线性二阶方程：
\begin{align}
\nabla\cdot\bigl(\rho(|\nabla\chi|^2,\chi,\rho_0)\nabla\chi\bigr)+2\rho(|\nabla\chi|^2,\chi,\rho_0)=0.
\label{align:ch18-18312}
\end{align}
回忆 $h'=c^2/\rho$，\eqref{align:ch18-18312}可写为
\begin{align}
c^2\Delta\chi-\sum_{i,j=1}^2\chi_{\xi_i\xi_j}\chi_{\xi_i}\chi_{\xi_j}=|\nabla\chi|^2-2c^2.
\label{align:ch18-18313}
\end{align}
以（伪）Mach 数
\begin{align}
L=\frac{|\nabla\chi|}{c}
\label{align:ch18-18314}
\end{align}
衡量，\eqref{align:ch18-18312}、\eqref{align:ch18-18313} 在 $L<1$ 的点处为椭圆型，在 $L>1$ 的点处为双曲型。

检验方程\eqref{align:ch18-18312}类型有一个富有启发性的几何方法如下。对任一固定状态 $(\rho,v)$，在 $\mathbb{R}^2$ 上配以圆心为 $v$、半径为 $c(\rho)$ 的圆 $|\xi-v|=c(\rho)$，称之为 $(\rho,v)$ 的声速圆。于是现在清楚：给定流动 $(\rho(\xi),v(\xi))$，\eqref{align:ch18-18312}在点 $\bar\xi$ 处为椭圆型（或双曲型），当且仅当 $\bar\xi$ 位于状态 $(\rho(\bar\xi),v(\bar\xi))$ 的声速圆内部（或外部）。

\eqref{align:ch18-18312}的弱解是满足该方程（在分布意义下）的 Lipschitz 函数 $\chi$。激波对应于 $\chi$ 的偏导数发生跳跃间断的曲线。由于 $\chi$ 本身连续，其切向导数跨过激波也必连续，故伪速度与速度的切向分量不能跳跃。特别地，（伪）流动不能支撑（伪）涡旋面。与之相对，$\chi$ 的法向导数、从而伪速度与速度的法向分量跨过激波必跳跃，其跳跃满足 Rankine--Hugoniot 条件
\begin{align}
\left[\!\left[\rho\bigl(|\nabla\chi|^2,\chi,\rho_0\bigr)\frac{\partial\chi}{\partial\nu}\right]\!\right]
=[\![\rho(v-\xi)\cdot\nu]\!]=0,
\label{align:ch18-18315}
\end{align}
其中 $\nu$ 为激波的单位法向量。Bernoulli 方程\eqref{align:ch18-1836}还给出额外的跳跃条件
\begin{align}
\left[\!\left[\frac12|\nabla\chi|^2+h(\rho)\right]\!\right]
=\left[\!\left[\frac12[(v-\xi)\cdot\nu]^2+h(\rho)\right]\!\right]=0.
\label{align:ch18-18316}
\end{align}

因而自相似流动与定常无旋流动的跳跃条件完全相同，只是前者中的伪速度 $u=v-\xi$ 扮演后者中速度的角色。特别地，激波极曲线在此同样适用。

应当注意，点 $\xi$ 处的激波速率为 $\xi\cdot v$。故直线激波以常速率传播。特别地，过原点的直线激波是定常的。此时 $\xi\cdot v=0$，\eqref{align:ch18-18315}、\eqref{align:ch18-18316}退化为定常流动的跳跃条件\eqref{align:ch18-18216}、\eqref{align:ch18-18217}。

与定常情形一样，任一可容激波的前面状态必为伪超声速，而后面状态可以是伪超声速、伪声速或伪亚声速；相应地，激波分别称为超声速、声速或跨声速激波。

我们期望：亚声速区域中\eqref{align:ch18-18313}的椭圆性所诱导的那一点稳定性，将阻止生成那种亚声速、声速与超声速区域精细交织的极复杂流型。下述关于\eqref{align:ch18-18313}解的极大值原理，为这一预期提供了佐证，其证明见\cref{section:ch18-注记}所引文献。

\begin{theorem}[椭圆性原理]\label{theorem:ch18-椭圆性原理}
设 $\chi$ 是\eqref{align:ch18-18313}在开有界集 $\Omega$ 上的光滑解，满足 $L\leqslant 1$ 且 $c\leqslant\bar c$。则存在依赖于 $\Omega$ 的 $\delta>0$，使得或者 $L^2\leqslant 1-\delta$；或者对任一满足 $|\nabla f|\leqslant\delta/\bar c$ 且 $|\nabla^2 f|\leqslant\delta/\bar c^2$ 的函数 $f\in C^2(\Omega)$，函数 $L^2+f$ 不在 $\Omega$ 内取到最大值。
\end{theorem}

由此推知：$\Omega$ 上在边界保持 $L\leqslant 1$ 的小扰动，也保持内部的椭圆性。\cref{theorem:ch18-椭圆性原理}的另一个推论是：不存在在开集上处处声速的光滑流动。

既然声速区域必为“薄”的，合理的猜测是：在一般流动中，亚声速与超声速区域之间的界面是一条曲线。它可以是一道跨声速激波，密度与速度跨过它跳跃，其跳跃满足\eqref{align:ch18-18315}、\eqref{align:ch18-18316}。或者，从亚声速到超声速的过渡可连续地发生，跨过一条声速曲线。为洞察后一情形，下面在最简单的设定下考察：一个亚声速区域与一个作为常状态 $(\bar\rho,\bar v)$ 的超声速区域相接。于是如前所述，界面是圆心在点 $\bar v$、半径为 $\bar c=c(\bar\rho)$ 的声速圆弧。不失一般性设 $\bar v=0$，以极坐标 $(r,\theta)$ 表示伪势 $\chi$：$\xi_1=r\cos\theta$，$\xi_2=r\sin\theta$。在超声速区域 $r>\bar c$，有 $\chi=\frac12 r^2$。我们关心 $\chi$ 在声速圆附近、亚声速一侧 $r<\bar c$ 的行为。$\rho$ 与 $v$ 跨过声速圆连续，故 $\chi$ 及其梯度在声速圆上连续。关键问题是：$\chi$ 的梯度跨过界面是否 Lipschitz、甚至更光滑？为确定起见，设气体为理想气体，绝热指数 $\gamma>1$，状态方程规范化为
\begin{align}
p(\rho)=\frac{\gamma-1}{\gamma}\rho^\gamma,\qquad c^2(\rho)=(\gamma-1)\rho^{\gamma-1},\qquad h(\rho)=\rho^{\gamma-1}.
\label{align:ch18-18317}
\end{align}
在极坐标下，方程\eqref{align:ch18-18313}取形式
\begin{align}
[c^2-(\chi_r-r)^2]\chi_{rr}-\frac{2}{r^2}(\chi_r-r)\chi_\theta\chi_{r\theta}
+\frac{1}{r^2}\left(c^2-\frac{1}{r^2}\chi_\theta^2\right)\chi_{\theta\theta}
+\frac{c^2}{r}\chi_r+\frac{1}{r^3}(\chi_r-2r)\chi_\theta^2=0.
\label{align:ch18-18318}
\end{align}

为仔细考察 $\chi$ 及其导数跨过界面的过渡，在\eqref{align:ch18-18318}中引入新变量 $s=\bar c-r$ 与 $w=\chi-\frac12 r^2$ 以代替 $r$ 与 $\chi$。注意到 $w,w_\theta,w_{\theta\theta},w_s,w_{s\theta},w_{s\theta\theta}$ 在 $s=0$ 处皆为零，比较\eqref{align:ch18-18318}中各项预期的数量级，可设：当 $s>0$ 充分小时，方程退化为
\begin{align}
[2s-(\gamma+1)w_s]w_{ss}-w_s=R,
\label{align:ch18-18319}
\end{align}
其中 $R=o(s)$。这一设定的证明以及关于解的精细估计，可在\cref{section:ch18-注记}所引文献中找到。此处我们将 $R=o(s)$ 视为当然，只勾勒
\begin{align}
w_s=\frac{s}{\gamma+1}+o(s)
\label{align:ch18-18320}
\end{align}
的推导；它表明 $(\rho,v)$ 跨过声速圆的过渡是 Lipschitz 的而非 $C^1$ 的，并且在 $s>0$ 充分小时流动为亚声速。

作变量替换
\begin{align}
w=\frac{z}{\gamma+1}+\frac{s^2}{\gamma+1},
\label{align:ch18-18321}
\end{align}
则\eqref{align:ch18-18319}给出
\begin{align}
z_s z_{ss}+3z_s+2s=o(s),
\label{align:ch18-18322}
\end{align}
它可积出
\begin{align}
z_s^2+6z+2s^2=o(s^2).
\label{align:ch18-18323}
\end{align}
注意当 $s>0$ 充分小时，由\eqref{align:ch18-18323}有 $z<0$，由\eqref{align:ch18-18322}有 $z_s\neq0$，故 $z_s<0$。变换 $z=-\frac16(\varpi^2+2)s^2$ 将\eqref{align:ch18-18323}化为关于 $\varpi$ 的可分离变量方程，易得 $\varpi=1+o(1)$，从而 $z=-\frac12 s^2+o(s^2)$，$z_s=-s+o(s)$，即证得\eqref{align:ch18-18320}。

注意到 $w=s^{3/2}$ 是线性方程 $2sw_{ss}-w_s=0$ 的解，便知\eqref{align:ch18-18319}中的非线性项 $w_sw_{ss}$ 在保证 $w_s$ 于 $s=0$ 处 Lipschitz 连续中起关键作用。

本节最后指出：Euler 方程组\eqref{align:ch03-等熵Euler}取 $b=0$，也容许自相似解 $\rho=\rho(\xi)$、$v=v(\xi)$，它们满足方程组
\begin{align}
\begin{cases}
-\xi\cdot\nabla\rho+\nabla\cdot(\rho v)=0\\
-\xi\cdot\nabla(\rho v)+\nabla\cdot(\rho v\otimes v)+\nabla p(\rho)=0,
\end{cases}
\label{align:ch18-18324}
\end{align}
或用伪速度 $u=v-\xi$ 表示，
\begin{align}
\begin{cases}
\nabla\cdot(\rho u)=-2\rho\\
\nabla\cdot(\rho u\otimes u)+\nabla p(\rho)=-3\rho u.
\end{cases}
\label{align:ch18-18325}
\end{align}
同样，自相似流动与定常非无旋流动的跳跃条件完全相同，前者中的伪速度扮演后者中速度的角色。

把 $\xi$ 坐标中 $(\rho,v)$ 的构型看作原始 $x$ 变量中实际流动在 $t=1$ 时刻的快照，往往是富有教益且有帮助的。

上述内容的应用将在接下来三节中出现。

\end{document}
